\scanpage{67}
\mspage{21}
\[
(23)\qquad I\xrightarrow{d}\mathbb Z
\qquad (\text{i.e. } I=\coprod_m I_m)
\]
jouant le rôle d'une fonction dimension (on aurait une présentation en quelque sorte duale avec une fonction type codimension), satisfaisant la condition
\[
(24)\qquad\text{Si }i,j\in I\text{ et }i\leq j,\text{ alors }d(j)=d(i)+1\Longleftrightarrow j\text{ succ. de }i
\]
Posons alors, pour $m,n_{\deleted{\unsure{2|q}}}\in\mathbb Z$
\[
(25)\qquad\left\{\begin{aligned}
X_m^*&=\coprod_{\substack{i\in I_m\\\deleted{\unsure{d(i)}}}}X_i^*\\[1.5em]
V_{m,\deleted{\unsure{n}}\inserted{m+1}}^{\deleted{\unsure{*}}}&=\coprod_{\substack{(i,j)\in I\times I\\i\leq j,\ d(i)=m\\d(j)=m+1}}V_{i,j}\\[1.5em]
V_{m,m+1}^*&=\coprod_{\substack{i,j\text{ comme}\\\text{dessus}}}V_{ij}^*
\end{aligned}\right.
\]
de sorte qu'on a les diagrammes
\[
\underset{(\mathrm{Diag})_m}{(26)}\qquad
\diagram{\begin{tikzcd}[column sep=large,row sep=large]
X_m^* \arrow[r] & V_{m,m+1}^{\deleted{\unsure{*}}} & V_{m,m+1}^* \arrow[l] \arrow[d] \\
& & X_{m+1}^*
\end{tikzcd}}
\]
s'insérant dans un diagramme commun \inserted{qui n'est autre que (20),} infini en principe, mais que nous écrirons sous forme finie en \deleted{\unsure{présumant|privilégiant}} \unsure{simplifiant|simplifier} \deleted{$I$ fini}, \inserted{$d(I)$ \unsure{fini et} $\subset\mathbb N$)} et en n'écrivant pas les \deleted{\unsure{composantes}} \inserted{parties du diagramme} correspondant à des sommets occupés par des espaces vides :

\mspage{66}
